% Resubmission manuscript, targeting Developmental Cognitive Neuroscience
% (Elsevier; free-format first submission, name-date citations). Link this
% repo to Overleaf (Menu -> GitHub) and compile with pdfLaTeX + BibTeX.
% Section drafts are tracked in docs/ until they stabilise; the Introduction
% below mirrors docs/INTRODUCTION_DRAFT.md.
\documentclass[11pt]{article}

\usepackage[utf8]{inputenc}
\usepackage[T1]{fontenc}
\usepackage[round]{natbib}
\usepackage{graphicx}
\usepackage{amsmath}
\usepackage[margin=2.5cm]{geometry}
\usepackage{lineno}
\usepackage{setspace}
\usepackage[hidelinks]{hyperref}

\title{The structural connectivity that distinguishes adult prefrontal
subregions is present in neonates\\
\large [title under discussion -- see docs/INTRODUCTION\_DRAFT.md for candidates]}

\author{Chiara Caldinelli \and Rhodri Cusack}
\date{}

\begin{document}
\maketitle
\linenumbers
\onehalfspacing

\begin{abstract}
[To be rewritten last, under the organisers framing: geometry specifies much
of adult DLPFC functional topography; connectivity carries a function-specific
contribution beyond geometry for all five task domains; that contribution is
already present in neonates.]
\end{abstract}

\section{Introduction}

At term birth, a human is still years away from comprehending a story, months
from reaching for an object, and more than a decade from the mature working
memory of an adult. Yet in every neurotypical adult, these functions occupy a
consistent arrangement of subregions within dorsolateral prefrontal cortex
(DLPFC): topographies of task preference that recur, individual after
individual, across working memory, language, social cognition, emotion and
action \citep{barch2013,glasser2016}. This consistency implies that something
specifies which patch of prefrontal cortex comes to serve which function, and
that the specifying information is in place by the time each function
emerges. Here we ask what that information is, and how early it can be
measured.

Three broad classes of answer are on offer. The first is \emph{geometry}: a
patch's fate follows from its position in the cortical sheet. Positional
identity is established early by morphogen gradients
\citep{rakic1988,oleary1989}, macroscale function is organised along smooth
spatial gradients \citep{margulies2016}, and much of adult functional
organisation is predictable from the geometry of the cortex alone
\citep{pang2023}. The second is \emph{extrinsic connectivity}: a patch
computes what its connections allow, so function follows the connectional
fingerprint \citep{passingham2002} --- an account with causal support from
rewiring experiments \citep{vonmelchner2000} and correlational support from
studies predicting adult functional selectivity from connectivity
\citep{saygin2011,osher2016}. The third is \emph{local microstructure}:
cytoarchitectonic identity could determine both a region's connections and
its eventual function, making connectivity a marker rather than a mechanism
\citep{kubota2025}. These accounts are not exclusive, and in any adult
dataset their signatures are correlated: connectivity is itself spatially
organised, and microstructure varies smoothly with position. A demonstration
that any one of them predicts function therefore does not, by itself,
distinguish between them.

The entanglement is not hypothetical. Diffusion tractography is biased toward
nearby targets, so a vertex's tractography profile partly encodes its
position; and group-average activation maps are spatially smooth, so a model
needs little more than position to approximate them. A connectivity-based
prediction of function can therefore succeed as a re-description of geometry
\citep{donahue2016,pang2023}. We address this by comparison against explicit
geometric baselines. We measure what spatial smoothness alone achieves, by
passing surrogate maps with matched spatial autocorrelation
\citep{burt2020} through the full analysis pipeline; what geometry achieves,
with a model given only each vertex's distances to the rest of the brain; and
what connectivity adds beyond both, including after the distance dependence
of tractography is removed from its profiles. We also estimate how the
variance of each functional map divides into geometric, connectional and
shared components (Fig.~1b).

These specifiers differ most sharply in \emph{when} their information becomes
available relative to experience. We make no claim that neonatal DLPFC is
inactive or uncommitted: it shows organised spontaneous activity, responses
to speech are detectable in frontal cortex even in premature neonates
\citep{mahmoudzadeh2013,dehaenelambertz2002}, attentional signals appear
within months \citep{ellis2021}, and neural activity can reveal function well
before behaviour does --- infants hold visual category information that overt
behaviour alone would not show \citep{odoherty2026}. Our premise is narrower
and behavioural: the functions indexed by our five contrasts --- instructed
working memory, narrative comprehension, explicit mentalising, cued action
selection, emotion recognition --- constitute an \emph{adult task repertoire}
that emerges over years and matures into the second decade
\citep{wellman2001,luna2004}, and DLPFC is among the last cortical
territories to mature by every anatomical index
\citep{chugani1986,huttenlocher1997,gogtay2004}. At term, essentially all of
the experience of performing this repertoire lies in the future. Whatever
precursor computations neonatal DLPFC performs, a specification of the
repertoire's adult topography that is measurable at birth cannot be a
consequence of performing that repertoire. Connectivity-based
proto-organisation of this kind has been demonstrated in ventral temporal
cortex, where connectivity at age five predicts where the visual word form
area will emerge by age eight \citep{saygin2016}. Here we ask the
corresponding question at the beginning of postnatal life, for five domains
at once, and under explicit geometric controls: is the \emph{connectional}
component of adult DLPFC topography --- the part geometry cannot carry ---
already present in the neonatal connectome?

We studied the topography of five task activations within DLPFC in 155 adults
\citep{vanessen2013,barch2013} and structural connectivity in the same adults
and in 325 term neonates \citep{edwards2022}, within a single predictive
framework: identical models, targets and metrics throughout, so that
competing accounts differ only in the information they are allowed. Four
predictions follow from the accounts above. (P1) If functional topography is
unconstrained smooth variation, real task maps should transfer from adults to
neonates no better than surrogate maps of matched smoothness. (P2) If
geometry suffices, distances alone should predict the maps as well as
tractography does, and removing distance information from tractography should
abolish its predictive power. (P3) If extrinsic connectivity contributes, it
should predict adult topography beyond these geometric baselines, with task
specificity. (P4) If that connectional contribution requires postnatal
experience of the repertoire, it should be absent at birth; if it is part of
the scaffold, it should be present. We find that geometry is a strong predictor
of these topographies --- for the motor map it outperforms tractography ---
but that connectivity carries a function-specific contribution beyond
geometry in all five domains, and that this contribution is already present
in the first weeks of life.

\section{Methods}
% DCN uses the conventional IMRaD order: Methods precede Results in the main
% flow (unlike the PNAS version, which had Materials and Methods at the end).
[Port from the submitted version with the corrections in
MANUSCRIPT\_CHANGES.md; add the control analyses. With Methods up front, the
Results need not carry capsule method descriptions --- define the ladder of
models here once.]

\section{Results}
[From PHASE0\_RESULTS.md and MANUSCRIPT\_CHANGES.md; to be drafted, walking
predictions P1--P4 in order.]

\section{Discussion}
[To be drafted; includes the microstructure third-variable discussion
(\citealp{kubota2025}) and the retirement of the region-level
``connectivity precedes function'' framing.]

\bibliographystyle{plainnat}
\bibliography{references}

\end{document}
